%% ============================================================
%%  Capítulo 5 – Descripción Técnica: PCB-PWR
%% ============================================================
\chapter{Descripción Técnica: PCB-PWR}
\label{chap:pcb_pwr}

\section{Función Principal}
\label{sec:pwr_funcion}

La PCB-PWR es el módulo de gestión de energía de ALMA. Su función es convertir la
entrada de 12\,V DC proveniente del adaptador externo en los rieles de baja tensión
regulados que alimentan todos los subsistemas de la PCB-CTRL. Adicionalmente,
concentra las protecciones de entrada del sistema: fusible PPTC, diodo TVS y
protección contra polaridad inversa. El subsistema alojado es el \textbf{SGE}
(Gestión de Energía) \reqref{REQ-E-01}.

\section{Bloques Funcionales Incluidos}
\label{sec:pwr_bloques}

\subsection{Etapa de Entrada y Protección}

La entrada de 12\,V VIN incorpora las siguientes protecciones, en el orden en que
aparecen en la señal de entrada:

\begin{itemize}
  \item \textbf{Protección de polaridad inversa:} diodo Schottky de baja caída en
    serie, o MOSFET-P como interruptor controlado, para prevenir daño por conexión
    inversa del adaptador.
  \item \textbf{Fusible PPTC resetable:} limita la corriente en caso de
    cortocircuito y se recupera al retirar la falla. Disparo en $<$5\,s según
    \reqref{REQ-E-01}.
  \item \textbf{Diodo TVS:} supresión de transitorios de tensión en la entrada,
    protegiendo todos los reguladores aguas abajo \reqref{REQ-E-01}.
  \item \textbf{Condensadores bulk de entrada:} capacidad suficiente para absorber
    los picos de corriente de los reguladores Buck durante transitorios de carga.
\end{itemize}

\subsection{Reguladores de Voltaje}

El árbol de potencia replica la topología del diseño de referencia Khadas VIM4,
adaptada a las necesidades de ALMA. La Tabla~\ref{tab:reguladores} resume los
reguladores implementados.

\begin{table}[H]
  \centering
  \caption{Reguladores de voltaje de la PCB-PWR.}
  \label{tab:reguladores}
  \begin{tabularx}{\textwidth}{L{2.2cm} C{2.5cm} C{1.8cm} C{2cm} L{4cm}}
    \toprule
    \rowcolor{udea-green}
    \textcolor{white}{\textbf{Componente}} &
    \textcolor{white}{\textbf{Fabricante}} &
    \textcolor{white}{\textbf{Topología}} &
    \textcolor{white}{\textbf{V salida}} &
    \textcolor{white}{\textbf{Destino / Función}} \\
    \midrule
    SY88003   & Silergy          & Buck síncrono & 1.1\,V   & VDD\_NPU: núcleo de IA del A311D2 \\
    \rowcolor{row-alt}
    TL76701   & Texas Instruments & LDO          & 1.8\,V   & VDD\_IO: bancos de E/S y buses lógicos \\
    MPU8756   & Monolithic Power  & Buck síncrono & Variable & VDD\_CPU: clústeres CPU-A y CPU-B \\
    \rowcolor{row-alt}
    NB679     & Monolithic Power  & Buck síncrono & 5.0\,V   & Amplificador Clase~D, backlight, LEDs \\
    NB680     & Monolithic Power  & Buck síncrono & 3.3\,V   & WiFi/BT, eMMC, GPIO, sensores \\
    \bottomrule
  \end{tabularx}
\end{table}

\subsection{Estrategia de Ajuste de Voltaje (Red de Feedback)}

Los reguladores Buck SY88003 y MPU8756 utilizan redes de resistencias de
\textit{feedback} (RFB) para definir con precisión el voltaje de salida. Esta es una
característica crítica del diseño:

\begin{itemize}
  \item El \textbf{SY88003} se emplea en múltiples instancias del esquemático
    (una por riel), cada una con valores de RFB calculados individualmente para
    entregar la tensión exacta requerida por cada dominio del A311D2. Por ejemplo,
    la instancia del riel NPU usa una red que entrega exactamente 1.1\,V, siguiendo
    los valores del esquemático de referencia de Khadas.
  \item El \textbf{MPU8756} gestiona los rieles variables de los clústeres CPU-A
    y CPU-B del A311D2 (arquitectura big.LITTLE). Sus voltajes se ajustan
    dinámicamente; los valores de RFB definen el rango de operación.
\end{itemize}

\begin{riskbox}[Mandato de Diseño: No Copiar-Pegar Componentes de Potencia]
  Cada instancia del SY88003 en el esquemático requiere valores de resistencias de
  feedback \textbf{calculados individualmente}. Copiar una instancia sin actualizar
  los valores pasivos resultaría en un voltaje de salida incorrecto, con riesgo de
  daño permanente al A311D2. Todos los valores de RFB deben verificarse contra la
  tabla de voltajes del esquemático de referencia Khadas VIM4 antes del cierre
  del esquemático.
\end{riskbox}

\subsection{Conector de Interconexión hacia PCB-CTRL}

Todos los rieles regulados se conducen desde la PCB-PWR hacia la PCB-CTRL a través
de un único conector de interconexión de alta corriente. Los pines del conector
se distribuyen de forma que los rieles de mayor corriente (VCC5 para el amplificador,
VDD\_CPU para los núcleos del SoC) dispongan de múltiples pines en paralelo para
reducir la resistencia de contacto y minimizar la caída de tensión.

\section{Interfaces de la PCB-PWR}
\label{sec:pwr_interfaces}

\begin{table}[H]
  \centering
  \caption{Interfaces de la PCB-PWR.}
  \label{tab:interfaces_pwr}
  \begin{tabularx}{\textwidth}{L{3cm} C{3cm} L{7cm}}
    \toprule
    \rowcolor{udea-green}
    \textcolor{white}{\textbf{Interfaz}} &
    \textcolor{white}{\textbf{Tipo}} &
    \textcolor{white}{\textbf{Descripción}} \\
    \midrule
    Entrada VIN    & Conector DC Jack o Molex & 12\,V DC desde adaptador externo \\
    \rowcolor{row-alt}
    Salida a CTRL  & Conector inter-PCB       & Rieles regulados: 1.1V, 1.8V, 5V, 3.3V, VDD\_CPU, GND \\
    \bottomrule
  \end{tabularx}
\end{table}

\section{Señales y Rieles Críticos}
\label{sec:pwr_senales}

\subsection{Rieles VDD\_CPU y VDD\_NPU (1.1\,V)}

Son los rieles más críticos de toda la plataforma. Una desviación de $\pm$5\,\%
puede provocar inestabilidad o daño al SoC. Requieren:
\begin{itemize}
  \item Condensadores de salida de baja ESR directamente en la salida del regulador.
  \item Pistas de potencia de ancho suficiente para la corriente máxima de los
    núcleos (calcular con IPC-2152, ver Anexo~\ref{chap:anexo_si}).
  \item Plano de cobre en capa interna dedicado a este riel, con múltiples vías de
    conexión al plano.
\end{itemize}

\subsection{Riel VCC5 -- Amplificador Clase~D}

El amplificador de audio Clase~D tiene corrientes pico elevadas durante la
reproducción a máxima potencia. Este riel debe contar con condensadores bulk
de alta capacitancia (electrolíticos + cerámicos) cerca del amplificador en la
PCB-CTRL para absorber los transitorios, evitando que se propaguen hacia atrás
hasta la PCB-PWR.

\subsection{Ruido de Switching}

Los reguladores Buck generan ruido de conmutación en frecuencias típicas de
200\,kHz a 2\,MHz. Este ruido es el motivo principal de la separación física de la
PCB-PWR. En la PCB-PWR misma, se deben seguir las siguientes buenas prácticas:
\begin{itemize}
  \item El lazo de alta corriente (inductor -- diodo/MOSFET -- condensador de salida)
    debe ser mínimo en área.
  \item Las pistas del nodo de conmutación (SW) deben ser cortas y no actuar
    como antenas.
  \item El plano de GND de la PCB-PWR debe ser continuo bajo los reguladores.
\end{itemize}

\section{Restricciones del Fabricante Aplicables a PCB-PWR}
\label{sec:pwr_fabricante}

La PCB-PWR puede fabricarse en un proceso estándar de \textbf{4~capas}. Las
restricciones relevantes del fabricante (documentadas en detalle en el
Capítulo~\ref{chap:fabricante}) son:

\begin{fabbox}[Restricciones Fabricante -- PCB-PWR]
  \begin{itemize}
    \item \textbf{Ancho mínimo de pista:} 0.1\,mm (aplicable a pistas de señal;
      las pistas de potencia de alta corriente requieren mayor ancho según IPC-2152).
    \item \textbf{Clearance:} 0.1\,mm general; distancias IEC\,62368-1 para la
      entrada de 12\,V \reqref{REQ-E-01}.
    \item \textbf{Acabado superficial:} ENIG recomendado para mejorar soldabilidad
      de componentes SMD de potencia.
    \item \textbf{Thermal relief en pads de componentes de potencia:} configurar
      según recomendación del fabricante para asegurar soldabilidad en reflow.
  \end{itemize}
\end{fabbox}

\section{Restricciones de Integridad de Señal -- PCB-PWR}
\label{sec:pwr_si}

Las señales en la PCB-PWR son exclusivamente DC y baja frecuencia de control.
Las restricciones de SI se centran en integridad de potencia (PDN):

\begin{sibox}[SI-PWR-01: PDN -- Estabilidad de Rieles Críticos]
  Los rieles VDD\_CPU y VDD\_NPU requieren una impedancia de PDN suficientemente baja
  para evitar caídas de tensión durante transitorios de carga del SoC.
  Criterio: variación máxima de $\pm$3\,\% en el riel bajo carga transitoria.
  Implementación: condensadores de salida de baja ESR en paralelo (cerámico 100\,nF
  + electrolítico 100\,$\mu$F mínimo en cada riel crítico).
\end{sibox}

\section{Captura en Altium Designer}
\label{sec:pwr_altium}

El esquemático de la PCB-PWR está estructurado por regulador individual, siguiendo
la organización del esquemático de referencia de Khadas VIM4. Cada regulador ocupa
una hoja esquemática independiente con su red de feedback, componentes pasivos de
filtrado y etiquetas de red (Net Labels) que coinciden con los nombres de riel de
la Tabla~\ref{tab:rieles_sistema}. Los pines del conector de interconexión
están etiquetados con los mismos nombres de red para garantizar la trazabilidad
entre las dos tarjetas en el entorno multi-board de Altium.
